% =====================================================================
% Extração de Entidades e Relações (Text-to-Graph)
% Responsável: Pessoa 2 (Anna Maria Albino)
%
% O que cobrir:
% - Como os textos foram divididos em chunks (tamanho, sobreposição) e por quê.
% - Qual LLM foi usado e com quais parâmetros (temperatura etc.).
% - O prompt de extração: como ele impõe os tipos do modelo (Tabelas 2 e 3) e as regras contra alucinação de docs/formato_extracao.md. Pode ir um trecho do prompt.
% - Formato de saída (JSON validado por src/contrato.py) e o que acontece com extrações inválidas.
% - Números: total de entidades e relações extraídas, % de JSONs aprovados pelo validador de primeira.
% - Avaliação manual numa amostra: quantas relações estavam corretas.
%
% Dicas: use \rot{Pessoa} para rótulos e \rel{CRIOU} para relações.
% Cite com \cite{chave} (chaves em referencias.bib). Ao terminar,
% apague o \pendente abaixo.
% =====================================================================

\section{Extração de Entidades e Relações (Text-to-Graph)}
\label{sec:extracao}

\pendente{extração de entidades e relações (text-to-graph)}
